\documentclass[12pt,a4paper]{article}
\usepackage[utf8]{vietnam}
\usepackage{amsmath,amssymb}
\usepackage{geometry}
\geometry{top=1.5cm,bottom=1.2cm,left=1.4cm,right=1.4cm,headheight=28pt,headsep=12pt,footskip=18pt}
\usepackage{tikz}
\usepackage{xcolor}
\usepackage{enumerate}
\usepackage{graphicx}
\usepackage[version=4]{mhchem}
\usepackage{fancyhdr}
\usepackage{ex_test}

\renewcommand{\baselinestretch}{1.0}
\setlength{\parskip}{1pt}

\definecolor{hfRed}{RGB}{211,47,47}
\definecolor{hfGreen}{RGB}{139,195,144}

\pagestyle{fancy}
\fancyhf{}
\fancyhead[L]{\small\itshape\color{hfRed} Tài liệu lưu hành nội bộ}
\fancyhead[R]{\small\itshape\color{hfRed} Thầy \textbf{TRẦN VĂN THẠNH} - 0777.470.803}
\renewcommand{\headrulewidth}{0.8pt}
\renewcommand{\headrule}{\hbox to\headwidth{\color{hfGreen}\leaders\hrule height \headrulewidth\hfill}}

\fancyfoot[L]{\small\itshape\color{hfRed} CS1: 6/15 Nguyễn Hoàng, P. Kim Long, TP Huế \quad|\quad CS2: 24 Đặng Thái Thân, TP Huế}
\fancyfoot[R]{\small\bfseries\color{hfRed} Trang \thepage}
\renewcommand{\footrulewidth}{0.8pt}
\renewcommand{\footrule}{\hbox to\headwidth{\color{hfGreen}\leaders\hrule height \footrulewidth\hfill}}

\fancypagestyle{firstpage}{
  \fancyhf{}
  \renewcommand{\headrulewidth}{0pt}
  \fancyfoot[L]{\small\itshape\color{hfRed} CS1: 6/15 Nguyễn Hoàng, P. Kim Long, TP Huế \quad|\quad CS2: 24 Đặng Thái Thân, TP Huế}
  \fancyfoot[R]{\small\bfseries\color{hfRed} Trang \thepage}
  \renewcommand{\footrulewidth}{0.8pt}
  \renewcommand{\footrule}{\hbox to\headwidth{\color{hfGreen}\leaders\hrule height \footrulewidth\hfill}}
}

\begin{document}
\thispagestyle{firstpage}

\noindent
\begin{minipage}[t]{0.45\textwidth}
	\centering\fontsize{10pt}{12pt}\selectfont
	\textbf{SỞ GIÁO DỤC VÀ ĐÀO TẠO THÀNH PHỐ HUẾ}\\
	\textbf{TRƯỜNG THPT HOÁ CHÂU}\\
	\textbf{ĐỀ CHÍNH THỨC}\\
	\textit{(Đề thi có 03 trang)}
\end{minipage}\hfill
\begin{minipage}[t]{0.53\textwidth}
	\centering\small
	\textbf{KIỂM TRA GIỮA HỌC KỲ I -- NĂM HỌC 2025 - 2026}\\
	\textbf{MÔN: HOÁ HỌC -- LỚP 11}\\
	\textit{Thời gian làm bài: 45 phút (không kể thời gian phát đề)}\\
	\fbox{\textbf{MÃ ĐỀ: 102}}
\end{minipage}

\smallskip
\noindent\hrulefill
\smallskip

\noindent\textbf{Họ, tên thí sinh:} \dotfill\ \textbf{Lớp:} \dotfill\ \textbf{Số báo danh:} \dotfill

\smallskip
\noindent\textit{(Cho biết nguyên tử khối: H = 1; C = 12; N = 14; O = 16; Na = 23; Al = 27; S = 32; Cl = 35,5; K = 39; Ca = 40; Fe = 56; Cu = 64; Ba = 137).}
\smallskip

\noindent
\textbf{\large A. PHẦN TRẮC NGHIỆM (7,0 điểm)}

\smallskip
\noindent
\textbf{PHẦN I (3,0 điểm): Câu trắc nghiệm nhiều phương án lựa chọn.}\\
\textit{Thí sinh trả lời từ câu 1 đến câu 12. Mỗi câu hỏi thí sinh chỉ chọn một phương án.}

\Opensolutionfile{ans}[ans-hoa11-ma102hs]

\begin{ex}
Cho phương trình: $\ce{CH3COOH + H2O <=> CH3COO- + H3O+}$. Trong phản ứng thuận, theo thuyết Brønsted -- Lowry chất nào là acid?
\choice
{$\ce{CH3COOH}$.}
{$\ce{CH3COO-}$.}
{$\ce{H3O+}$.}
{$\ce{H2O}$.}
\loigiai{}
\end{ex}

\begin{ex}
Vai trò của $\ce{NH3}$ trong phản ứng: $\ce{4NH3 + 5O2 ->[t^\circ, \text{Pt}] 4NO + 6H2O}$ là
\choice
{chất khử.}
{chất oxi hóa.}
{acid.}
{base.}
\loigiai{}
\end{ex}

\begin{ex}
Để xác định nồng độ của một dung dịch $\ce{NaOH}$, người ta đã tiến hành chuẩn độ bằng dung dịch $\ce{HCl}\ 0{,}1\text{ M}$. Để chuẩn độ $10\text{ mL}$ dung dịch $\ce{NaOH}$ này cần $12\text{ mL}$ dung dịch $\ce{HCl}$. Nồng độ của dung dịch $\ce{NaOH}$ trên là
\choice
{$1{,}2\text{ M}$.}
{$0{,}012\text{ M}$.}
{$0{,}1\text{ M}$.}
{$0{,}12\text{ M}$.}
\loigiai{}
\end{ex}

\begin{ex}
Phản ứng thuận nghịch là phản ứng
\choice
{xảy ra giữa hai chất khí.}
{xảy ra theo hai chiều ngược nhau trong cùng điều kiện.}
{có phương trình hoá học được biểu diễn bằng mũi tên một chiều.}
{chỉ xảy ra theo một chiều nhất định.}
\loigiai{}
\end{ex}

\begin{ex}
$\ce{HNO3}$ tinh khiết là chất lỏng không màu, nhưng dung dịch $\ce{HNO3}$ để lâu thường chuyển sang màu vàng là do
\choice
{khi để lâu $\ce{HNO3}$ bị khử bởi các chất của môi trường.}
{$\ce{HNO3}$ tan nhiều trong nước.}
{dung dịch $\ce{HNO3}$ có tính oxi hóa mạnh.}
{dung dịch $\ce{HNO3}$ có hoà tan một lượng nhỏ $\ce{NO2}$.}
\loigiai{}
\end{ex}

\begin{ex}
Cho cân bằng (trong bình kín) sau:
\[ \ce{CO(g) + H2O(g) <=> CO2(g) + H2(g)} \quad \Delta_r H_{298}^\circ < 0 \]
Trong các yếu tố: (1) tăng nhiệt độ; (2) thêm một lượng hơi nước; (3) thêm một lượng $\ce{H2}$; (4) tăng áp suất chung của hệ; (5) dùng chất xúc tác. Dãy gồm các yếu tố đều làm thay đổi cân bằng của hệ là
\choice
{(2), (3), (4).}
{(1), (2), (3).}
{(1), (4), (5).}
{(1), (2), (4).}
\loigiai{}
\end{ex}

\begin{ex}
Phát biểu nào dưới đây \textbf{không đúng}?
\choice
{$\ce{H2SO4}$ đặc là chất hút nước mạnh.}
{Khi pha loãng sulfuric acid chỉ được cho từ từ nước vào acid.}
{$\ce{H2SO4}$ loãng có đầy đủ tính chất chung của acid.}
{Khi tiếp xúc với $\ce{H2SO4}$ đặc dễ gây bỏng nặng.}
\loigiai{}
\end{ex}

\begin{ex}
Cho các phát biểu sau:\\
(a) Sulfur là chất rắn, màu trắng, không tan trong nước.\\
(b) Sulfur và sulfur dioxide vừa có tính oxi hóa, vừa có tính khử.\\
(c) Khi phản ứng với hydrogen, sulfur thể hiện tính oxi hóa.\\
(d) Nước thải sinh hoạt là một trong các nguồn phát thải khí $\ce{SO2}$.\\
(e) Sulfur dioxide được sử dụng để tẩy trắng vải sợi, bột giấy, sản xuất sulfuric acid và diệt nấm mốc.\\
Số phát biểu \textbf{đúng} là
\choice
{4.}
{1.}
{2.}
{3.}
\loigiai{}
\end{ex}

\begin{ex}
Để tạo độ xốp cho một số loại bánh, có thể dùng chất nào sau đây?
\choice
{$\ce{NH4HCO3}$.}
{$\ce{CaCO3}$.}
{$\ce{NaCl}$.}
{$\ce{(NH4)3PO4}$.}
\loigiai{}
\end{ex}

\begin{ex}
Ở nhiệt độ thường, nitrogen khá trơ về mặt hoạt động là do
\choice
{Phân tử nitrogen có liên kết ba khá bền.}
{Nitrogen có độ âm điện lớn nhất trong nhóm.}
{Phân tử nitrogen không phân cực.}
{Nitrogen có bán kính nguyên tử nhỏ.}
\loigiai{}
\end{ex}

\begin{ex}
Tính chất nào sau đây \textbf{không phải} tính chất vật lí của sulfur?
\choice
{Thể rắn ở điều kiện thường.}
{Không tan trong nước.}
{Màu vàng ở điều kiện thường.}
{Không tan trong benzene.}
\loigiai{}
\end{ex}

\begin{ex}
Nitric acid thường được sử dụng để tẩy rửa ống, bề mặt kim loại do có
\choice
{tính khử.}
{tính oxi hoá mạnh.}
{tính acid yếu.}
{tính khử và tính acid mạnh.}
\loigiai{}
\end{ex}

\Closesolutionfile{ans}

\bigskip
\noindent
\textbf{PHẦN II (2,0 điểm): Câu trắc nghiệm đúng sai.}\\
\textit{Thí sinh trả lời từ câu 1 đến câu 2. Trong mỗi ý a), b), c), d) ở mỗi câu, thí sinh chọn đúng hoặc sai.}

\begin{ex}
Cho phản ứng tổng hợp ammonia:
\[ \ce{N2(g) + 3H2(g) <=>[xt, t^\circ, p] 2NH3(g)} \quad \Delta_r H_{298}^\circ = -92\text{ kJ} \]
\choiceTF
{a) Khi tăng nhiệt độ của hệ thì cân bằng chuyển dịch theo chiều nghịch.}
{b) Phản ứng trên là phản ứng thu nhiệt theo chiều thuận.}
{c) Để phản ứng tổng hợp $\ce{NH3}$ đạt hiệu suất cao thì người ta thực hiện ở nhiệt độ rất thấp và áp suất cao.}
{d) Trong phản ứng trên $\ce{N2}$ đóng vai trò là chất khử.}
\loigiai{}
\end{ex}

\begin{ex}
Chọn phát biểu đúng -- sai:
\choiceTF
{a) Dung dịch sulfuric acid đặc hòa tan được tất cả các kim loại.}
{b) Hòa tan $18{,}4\text{ gam}$ hỗn hợp X gồm Fe và Cu trong dung dịch $\ce{H2SO4}$ (đặc, nóng, dư), thu được $8{,}6765\text{ Lít}$ $\ce{SO2}$ (đkc) và dung dịch Y. Phần trăm khối lượng của Fe trong hỗn hợp X là $30{,}43\%$.}
{c) Dung dịch $\ce{HNO3}$ đặc, nguội phản ứng được với Al và Fe.}
{d) Sulfuric acid tan tốt trong nước, quá trình hòa tan tỏa nhiệt mạnh.}
\loigiai{}
\end{ex}

\bigskip
\noindent
\textbf{PHẦN III (2,0 điểm): Câu trắc nghiệm yêu cầu trả lời ngắn.}\\
\textit{Thí sinh trả lời từ câu 1 đến câu 8.}

\begin{ex}
Cho các phản ứng:\\
(a) $\ce{S + O2 ->[t^\circ] SO2}$\\
(b) $\ce{S + 3F2 -> SF6}$\\
(c) $\ce{Hg + S -> HgS}$\\
(d) $\ce{H2 + S ->[t^\circ] H2S}$\\
Số phản ứng trong đó sulfur đơn chất đóng vai trò chất khử là bao nhiêu?
\shortans{}
\loigiai{}
\end{ex}

\begin{ex}
Cho các chất: (1) $\ce{KNO3}$, (2) $\ce{HCl}$, (3) $\ce{HClO}$, (4) $\ce{H2S}$, (5) $\ce{Cu(OH)2}$, (6) $\ce{Ba(OH)2}$, (7) $\ce{NaHCO3}$. Có bao nhiêu chất điện li mạnh trong các chất trên?
\shortans{}
\loigiai{}
\end{ex}

\begin{ex}
Cho các phản ứng sau:\\
(1) $\ce{NaOH + HCl -> NaCl + H2O}$\\
(2) $\ce{H2 + I2 <=> 2HI}$\\
(3) $\ce{CaCO3 <=> CaO + CO2}$ (trong bình kín)\\
(4) $\ce{2KClO3 ->[t^\circ] 2KCl + 3O2}$\\
Có bao nhiêu phản ứng thuận nghịch trong các phản ứng trên?
\shortans{}
\loigiai{}
\end{ex}

\begin{ex}
Cho các nhận định sau về cấu tạo phân tử nitric acid:\\
(1) Liên kết $\ce{O - H}$ phân cực về oxygen.\\
(2) Nguyên tử N có số oxi hóa là $+5$.\\
(3) Nguyên tử N có hóa trị bằng 4.\\
(4) Phân tử $\ce{HNO3}$ chứa liên kết cho -- nhận $\ce{N -> O}$.\\
Có bao nhiêu phát biểu đúng?
\shortans{}
\loigiai{}
\end{ex}

\begin{ex}
Cho phản ứng: $a\ce{Fe} + b\ce{HNO3 ->} c\ce{Fe(NO3)3} + d\ce{NO ^} + e\ce{H2O}$. Hệ số tỉ lượng $a, b, c, d, e$ là những số nguyên dương có tỉ lệ tối giản. Tổng $(a + b)$ bằng bao nhiêu?
\shortans{}
\loigiai{}
\end{ex}

\begin{ex}
Cho sơ đồ chuyển hóa nitrogen trong khí quyển thành phân đạm:\\
$\ce{N2 ->[(1)][+O2] NO ->[(2)][+O2] NO2 ->[(3)][+O2, +H2O] HNO3 ->[(4)] NO3-}$.\\
Có bao nhiêu phản ứng thuộc loại oxi hóa -- khử trong sơ đồ trên?
\shortans{}
\loigiai{}
\end{ex}

\begin{ex}
Cho $m$ gam Fe tác dụng hết với dung dịch $\ce{HNO3}$ loãng (dư), sau phản ứng thu được $4{,}958\text{ lít}$ (đkc) khí NO là sản phẩm khử duy nhất. Tính giá trị của $m$?
\shortans{}
\loigiai{}
\end{ex}

\begin{ex}
Có các thí nghiệm sau:\\
(a) Nhúng thanh sắt (iron) vào dung dịch $\ce{H2SO4}$ loãng, nguội.\\
(b) Sục khí $\ce{SO2}$ vào nước bromine.\\
(c) Cho dung dịch $\ce{BaCl2}$ tác dụng với dung dịch $\ce{H2SO4}$ loãng.\\
(d) Nhúng lá nhôm (aluminium) vào dung dịch $\ce{H2SO4}$ đặc, nguội.\\
Số thí nghiệm xảy ra phản ứng hoá học là bao nhiêu?
\shortans{}
\loigiai{}
\end{ex}

\bigskip
\noindent
\textbf{\large B. PHẦN TỰ LUẬN (3,0 điểm)}

\begin{ex}
\textbf{(0,5 điểm)}
\begin{enumerate}[a)]
	\item Ở các vùng quê, người dân thường dùng phèn chua ($\ce{K2SO4.Al2(SO4)3.24H2O}$) để làm trong nước nhờ ứng dụng của phản ứng thuỷ phân ion $\ce{Al^3+}$. Giải thích?
	\item Chất hay ion nào là acid, là base trong phản ứng thuỷ phân $\ce{Al^3+}$?
\end{enumerate}
\loigiai{}
\vspace*{2.5cm}
\end{ex}

\begin{ex}
\textbf{(1,5 điểm)}
\begin{enumerate}[a)]
	\item Viết sơ đồ sản xuất nitric acid trong công nghiệp.
	\item Tính số kg dung dịch $\ce{HNO3}\ 60\%$ điều chế được từ $340\text{ kg}$ ammonia, biết rằng hiệu suất của toàn bộ quá trình là $90\%$.
\end{enumerate}
\loigiai{}
\vspace*{4.5cm}
\end{ex}

\begin{ex}
\textbf{(1,0 điểm)}
\begin{enumerate}[a)]
	\item Bình đựng dung dịch $\ce{H2SO4}$ đặc để trong không khí ẩm lâu ngày thì khối lượng bình có thay đổi không? Vì sao?
	\item Sulfur và quặng pyrite sắt là các nguyên liệu chính trong công nghiệp sản xuất sulfuric acid. Tại một nhà máy, cứ đốt cháy 1 tấn quặng pyrite sắt (chứa $84\%$ khối lượng $\ce{FeS2}$) bằng không khí, thu được tối đa $V\text{ m}^3$ khí $\ce{SO2}$ (đkc). Tính giá trị của $V$?
\end{enumerate}
\loigiai{}
\vspace*{4.5cm}
\end{ex}

\begin{center}
\textbf{--------- HẾT ---------}\\
\textit{(Học sinh không được sử dụng tài liệu. Giám thị không giải thích gì thêm)}
\end{center}

\end{document}
